\documentclass[12pt,a4paper]{article}
\usepackage[utf8]{inputenc}
\usepackage{amsmath,amssymb,amsthm,mathrsfs}
\usepackage[margin=2.5cm]{geometry}
\usepackage{hyperref}
\usepackage{enumitem}

\newtheorem{theorem}{Theorem}[section]
\newtheorem{proposition}[theorem]{Proposition}
\newtheorem{lemma}[theorem]{Lemma}
\newtheorem{corollary}[theorem]{Corollary}
\newtheorem{conjecture}[theorem]{Conjecture}
\theoremstyle{definition}
\newtheorem{definition}[theorem]{Definition}
\newtheorem{remark}[theorem]{Remark}
\newtheorem{question}[theorem]{Question}

\newcommand{\CC}{\mathbb{C}}
\newcommand{\ZZ}{\mathbb{Z}}
\newcommand{\NN}{\mathbb{N}}
\newcommand{\kk}{k}
\newcommand{\pp}{\mathfrak{m}}
\newcommand{\Aut}{\operatorname{Aut}}
\newcommand{\End}{\operatorname{End}}
\newcommand{\GJac}{G^{\mathrm{Jac}}}
\newcommand{\Ga}{\mathbb{G}_a}
\newcommand{\Gm}{\mathbb{G}_m}

\title{Synthesis of All Approaches to the\\
Jacobian Conjecture in Dimension Two\\[6pt]
\large Post-2026 Status Report}
\author{Research Notes --- October 2026}
\date{}

\begin{document}
\maketitle

\begin{abstract}
We synthesize the results of four parallel investigations into the Jacobian
Conjecture in dimension two, the last open case after Alp\"oge's July 2026
counterexample in dimension $\geq 3$. Each strategy is analyzed for its
strengths, known obstructions, and connections to the monomial algebra framework.
A potentially decisive development is identified: Zheglov's claimed proof
(October 2024) of the Dixmier Conjecture for $A_1$, which would settle the
problem via the Belov-Kanel--Kontsevich equivalence.
\end{abstract}

\tableofcontents

%% ====================================================================
\section{Overview: The Four Strategies}
%% ====================================================================

\begin{center}
\renewcommand{\arraystretch}{1.4}
\begin{tabular}{|c|l|l|}
\hline
\textbf{Strategy} & \textbf{Core Idea} & \textbf{Status} \\
\hline
A & Uniform degree bound for Anick approximants & Gap identified \\
B & Dixmier--Kontsevich equivalence (DC$_1 \Leftrightarrow$ JC$_2$) & \textbf{Claimed proof} \\
C & Properness via non-proper value set $S_F = \emptyset$ & Severe constraints known \\
D & Newton polygon + Demazure root compatibility & Degree $\leq 108$ done \\
\hline
\end{tabular}
\end{center}

%% ====================================================================
\section{Strategy A: Inverse Limit and Degree Bounds}
%% ====================================================================

\subsection{The Strategy}
Reduce JC$_2$ to showing that the Anick approximants $\varphi_d \in \Aut_\kk(\kk[x,y])$
satisfying $F \equiv \varphi_d \pmod{\pp^d}$ have $\deg \varphi_d \leq C(\deg F)$
independent of $d$.

\subsection{Known Results}
\begin{itemize}
\item \textbf{Gabber bound:} For an automorphism $F$ of $\kk[x,y]$,
  $\deg F^{-1} \leq \deg F$. This holds \emph{after} invertibility is established.
\item \textbf{Jung--van der Kulk:} $\Aut_\kk(\kk[x,y]) = \mathrm{Aff}_2 *_C \mathrm{Tri}_2$.
  Every automorphism has a reduced decomposition whose factor degrees are determined
  by the leading form.
\item \textbf{Furter:} The ``length'' function (number of factors in the JvdK
  decomposition) is lower semicontinuous. The variety $A_n$ of automorphisms of
  degree $\leq n$ has dimension $n + 6$ for $n \geq 2$ and is irreducible only
  for $n \leq 3$.
\item \textbf{Anick's theorem:} Gives no control over $\deg \varphi_d$. The proof
  is purely topological (density in the $\pp$-adic topology).
\end{itemize}

\subsection{Gap and Assessment}
No existing result gives a uniform degree bound on $\varphi_d$. Furter's
irreducibility failure for $A_n$ ($n \geq 4$) means the geometry of automorphism
varieties is complex, and families can exhibit unbounded degree growth while
converging $\pp$-adically.

\textbf{New input needed:} The absence of inner roots in dimension 2
(our Proposition from the main document) means the structure of $\GJac(A_d)$
is \emph{exactly} the image of tame automorphisms. Can this be used to control
degree via the amalgamated product structure? The key would be showing that
the ``reduced word length'' in the JvdK decomposition is bounded.

\textbf{Connection:} Karasek et al.\ showed that certain multidegrees in
$\Aut(\CC^3)$ would imply failure of JC$_2$, linking the problem to multidegree
theory.

%% ====================================================================
\section{Strategy B: The Dixmier--Kontsevich Path}
%% ====================================================================

\subsection{The Equivalence Chain}
The following equivalences are established:
\begin{enumerate}
\item \textbf{JC$_{2n} \Leftrightarrow$ DC$_n$:} Jacobian Conjecture in $2n$
  variables $\Leftrightarrow$ Dixmier Conjecture for $W_n$ (Belov-Kanel--Kontsevich
  2005, Tsuchimoto 2005).
\item \textbf{JC$_{2n} \Leftrightarrow$ PC$_n$:} $\Leftrightarrow$ Poisson Conjecture
  in $2n$ variables.
\item In particular: $\boxed{\text{JC}_2 \Leftrightarrow \text{DC}_1}$.
\end{enumerate}

After the 2026 counterexample:
\begin{itemize}
\item JC$_n$ is \textbf{false} for $n \geq 3$.
\item DC$_n$ is \textbf{false} for $n \geq 2$ (since JC$_4$ false $\Rightarrow$ DC$_2$ false).
\item PC$_n$ is \textbf{false} for $n \geq 2$.
\item \textbf{JC$_2$, DC$_1$, PC$_1$ remain open} (and equivalent).
\end{itemize}

\subsection{Zheglov's Claimed Proof of DC$_1$}

\begin{theorem}[Zheglov, arXiv:2410.06959, October 2024]
Every algebra endomorphism of the first Weyl algebra $A_1 = \CC\langle x, \partial_x \rangle
/ ([\partial_x, x] = 1)$ is an automorphism.
\end{theorem}

\textbf{If correct, this settles JC$_2$.}

\begin{itemize}
\item The proof uses the theory of normal forms for ordinary differential operators
  (generalized Schur theory) and the string equation.
\item The paper has undergone 4 revisions (25 pages $\to$ 67 pages).
\item Submitted to a journal on January 19, 2026.
\item As of October 2026: \textbf{no peer-reviewed publication; no known refutation;
  no community consensus.}
\end{itemize}

\subsection{The Kontsevich Isomorphism}

For $n = 1$, Kontsevich's conjecture that $\Aut(A_1) \cong \Aut_{\mathrm{symp}}(\CC^2)$
is \emph{confirmed}. The path is:
\[
  \text{Keller map } F: \CC^2 \to \CC^2 \xrightarrow{\text{symplectomorphism}}
  \text{End}(A_1) \xrightarrow{\text{DC}_1} \text{Aut}(A_1)
  \xrightarrow{\text{Kontsevich}} \Aut_{\mathrm{symp}}(\CC^2) = \Aut(\CC^2)
\]

\subsection{Assessment}
This is the \textbf{most promising path} to JC$_2$. The entire chain is
established except for DC$_1$ itself. The proof of DC$_1$ would be
purely algebraic (Weyl algebra theory), bypassing all geometric difficulties.
However, Zheglov's proof is unverified.

\textbf{Connection to our framework:} The Weyl algebra $A_1$ and the algebras
$A_d = \kk[x,y]/\pp^d$ are related via deformation quantization. The Demazure root
structure of $\Aut(A_d)$ should have a ``quantum'' counterpart in $\End(A_1)$.
Exploring this connection could provide an independent path to DC$_1$.

%% ====================================================================
\section{Strategy C: Properness and the Non-Proper Value Set}
%% ====================================================================

\subsection{The Framework}
A Keller map $F: \CC^2 \to \CC^2$ is \'etale, so failure of bijectivity
$\Leftrightarrow$ failure of properness $\Leftrightarrow$ $S_F \neq \emptyset$.

\subsection{Known Constraints on $S_F$}

\begin{enumerate}
\item \textbf{Jelonek (1993):} $S_F$ is either empty or a curve (union of algebraic curves).
\item \textbf{Nguyen Van Chau:} $S_F$ has exactly one point at infinity. Each irreducible
  component is parametrized by $(p(t), q(t))$ with $\deg p / \deg q = \deg P / \deg Q$.
\item \textbf{Nollet--Xavier:} $\pi_1(\CC^2 \setminus S_F)$ is non-abelian.
  $S_F$ cannot be a connected smooth hypersurface.
\item \textbf{Nollet--Taylor--Xavier:} $S_F$ cannot be a connected non-bifurcated
  algebraic hypersurface.
\item \textbf{Resolution at infinity:} Extending $F$ to $\PP^2 \dashrightarrow \PP^2$
  and resolving gives a weighted tree of exceptional curves. Dicritical divisors
  (Abhyankar) control non-properness. The constant Jacobian forces severe constraints
  on the resolution graph.
\end{enumerate}

\subsection{The Real Case (Pinchuk)}
The real Jacobian conjecture is \textbf{false}: Pinchuk (1994) gave a polynomial map
$\RR^2 \to \RR^2$ with everywhere-positive Jacobian that is not injective. Crucially,
the Jacobian is \emph{non-constant} (infimum is zero), so the map is non-proper.
This mechanism fails over $\CC$ because the constant Jacobian condition prevents
the Jacobian from approaching zero at infinity.

\subsection{Assessment}
The accumulated constraints make $S_F \neq \emptyset$ very difficult. The approach is:
show that for a Keller map, the constraints on $S_F$ are mutually contradictory.

\textbf{Connection to our framework:} The properness condition can be studied via the
inverse system $\{A_d\}$: $F$ is proper iff the induced maps $F_d: A_d \to A_d$ are
``uniformly surjective'' in a suitable sense. The outer root structure of $\GJac(A_d)$
forces each $F_d$ to be an automorphism, and the automorphisms of $A_d$ are proper
(finite-dimensional algebra). The question is whether this ``level-wise properness''
lifts to the polynomial level.

%% ====================================================================
\section{Strategy D: Newton Polygons and Root Compatibility}
%% ====================================================================

\subsection{Known Constraints on Counterexamples}

For a hypothetical counterexample $(P, Q)$ with $\det J(F) = 1$:

\begin{enumerate}
\item $\gcd(\deg P, \deg Q) \geq 16$.
\item $\gcd(\deg P, \deg Q) \neq p$ (prime) and $\neq 2p$.
\item Verified for $\max(\deg P, \deg Q) \leq 108$ (Guccione--Guccione--Horruitiner--Valqui).
\item The only surviving pair below max degree 125 is $(\deg P, \deg Q) = (72, 108)$
  or its swap.
\item Newton polygon of $P$ is a triangle with vertices $(0,0)$, $(n,0)$, $(0,m)$.
\item Top-degree edge factors as $c(x^\ell - y^k)^e$.
\item No edge parallel to the diagonal (slope 1).
\item Lee--Li (2024): the innermost vertex (NE vertex of the inner polynomial's
  Newton polygon) lies in a specific constrained region.
\item Semigroup constraints from Abhyankar--Moh theory via approximate roots and
  $G$-adic expansions.
\end{enumerate}

\subsection{Root-Polygon Compatibility}
The outer Demazure roots of $\pp^d$ in dimension 2 are:
\[
  R_1 = \{(-1, j) : j \geq 0\}, \quad R_2 = \{(i, -1) : i \geq 0\}.
\]
Each root direction corresponds to an edge direction of the Newton polygon.
The constraints above severely restrict which edge directions are compatible.

\subsection{Assessment}
The degree constraints are increasingly tight. The approach is inherently finite
(induction on degree), but the base case keeps growing. Current limit: degree 108.

\textbf{Key question:} Can the Demazure root framework give a \emph{structural}
(non-computational) reason why no valid Newton polygon exists for a non-invertible
Keller pair?

%% ====================================================================
\section{Interconnections and a Unified Strategy}
%% ====================================================================

The four strategies are not independent. We identify the following connections:

\begin{enumerate}
\item \textbf{A $\leftrightarrow$ D:} The Newton polygon determines the leading form
  of the JvdK decomposition. A structural result showing all edges are compatible with
  the outer root structure would give the degree bound of Strategy A.

\item \textbf{A $\leftrightarrow$ C:} A uniform degree bound on $\varphi_d$ implies
  $F = \varphi_d$ for large $d$ (an automorphism), which implies properness.
  Conversely, properness of $F$ combined with the Keller condition gives bijectivity
  (by standard theory), which gives the polynomial inverse.

\item \textbf{B $\leftrightarrow$ all:} The Dixmier path (Strategy B) is independent
  and would settle everything. But understanding \emph{why} DC$_1$ is true would connect
  back to the geometric and combinatorial structures of the other strategies.

\item \textbf{Monomial algebras as the bridge:} The key observation that $R_{\mathrm{inn}}(\pp^d)
  = \emptyset$ in dimension 2 is the combinatorial shadow of the algebraic fact that
  $A_1$ has no ``inner'' structure (it is generated by $x$ and $\partial_x$ with a
  single relation). This rigidity is precisely what makes DC$_1$ plausible (and true,
  if Zheglov is correct), while DC$_n$ fails for $n \geq 2$ where ``inner'' phenomena
  (wild automorphisms, tangent sweeps) appear.
\end{enumerate}

%% ====================================================================
\section{Proposed Research Program}
%% ====================================================================

\subsection{Immediate Priority: Verify Zheglov}

If Zheglov's proof of DC$_1$ is correct, JC$_2$ is settled. The first priority is:
\begin{itemize}
\item Read and verify Zheglov's proof (67 pages, arXiv:2410.06959v4).
\item Focus on the key steps: the theory of normal forms for differential operators,
  the string equation, and the passage from formal to algebraic statements.
\item Cross-check with the known constraints from Strategies A, C, D.
\end{itemize}

\subsection{Independent Paths (if Zheglov fails)}

\begin{enumerate}
\item \textbf{Root-polygon impossibility:} Show that the constraints from Strategy D,
  combined with the Demazure root structure, leave no valid Newton polygon for a
  counterexample. This would require extending the degree bound beyond 108 or finding
  a structural argument.

\item \textbf{Level-wise properness:} Develop a notion of ``uniform properness across
  truncation levels'' and show it implies properness of the polynomial map.

\item \textbf{Amalgamated product degree control:} Use the lower semicontinuity of
  the JvdK length (Furter) combined with the absence of inner roots to prove the
  uniform degree bound.

\item \textbf{Deformation quantization path:} Connect the Demazure root structure
  of $\Aut(A_d)$ to the endomorphism structure of $A_1$ via deformation quantization
  of the Poisson algebra $\CC[x,y]$, providing an independent proof of DC$_1$.
\end{enumerate}

\subsection{Computational Support}

\begin{itemize}
\item Extend degree verification beyond 108 using Newton polygon + root compatibility.
\item Implement Zheglov's normal form algorithm for small-degree differential operators.
\item Compute the ``quantum Demazure roots'' of the Weyl algebra truncations
  $A_1 / \mathfrak{m}^d_{A_1}$.
\item Systematic fiber computation for families of Keller maps of increasing degree.
\end{itemize}

%% ====================================================================
\section{Conclusion}
%% ====================================================================

The Jacobian Conjecture in dimension 2 appears to be on the verge of resolution.
The most direct path is through the Dixmier Conjecture for $A_1$ (Zheglov's claimed
proof). Even if this proof has gaps, the structural understanding from the monomial
algebra framework---particularly the absence of inner roots and the exact
correspondence between $\GJac(A_d)$ and tame automorphisms---provides multiple
independent avenues of attack.

The 2026 counterexample in dimension $\geq 3$ does not threaten the two-dimensional
case; rather, it \emph{clarifies} the situation by confirming that dimension 2 is
structurally different (all automorphisms tame, no inner roots, no tangent sweep
mechanism). The question is whether this structural rigidity can be turned into a proof.

\begin{thebibliography}{99}
\bibitem{Alpoge} L.~Alp\"oge, \emph{Counterexample to the Jacobian conjecture in
dimension 3}, July 2026.
\bibitem{Gao} S.~Gao, \emph{Counterexamples to the Jacobian conjecture in dimensions
$> 2$}, arXiv:2608.00222.
\bibitem{Zheglov} A.~Zheglov, \emph{The Conjecture of Dixmier for the first Weyl
algebra is true}, arXiv:2410.06959.
\bibitem{BKK} A.~Belov-Kanel, M.~Kontsevich, \emph{The Jacobian conjecture is stably
equivalent to the Dixmier conjecture}, Moscow Math.\ J.\ \textbf{7} (2007), 209--218.
\bibitem{Tsuchimoto} Y.~Tsuchimoto, \emph{Endomorphisms of Weyl algebra and
$p$-curvatures}, Osaka J.\ Math.\ \textbf{42} (2005), 435--452.
\bibitem{Jelonek} Z.~Jelonek, \emph{The set of points at which a polynomial map is
not proper}, Ann.\ Polon.\ Math.\ \textbf{58} (1993), 259--266.
\bibitem{NguyenVC} N.~V.~Chau, \emph{Non-proper value set and Jacobian condition},
preprint.
\bibitem{NolletXavier} S.~Nollet, F.~Xavier, \emph{A note on the Jacobian conjecture},
arXiv:2011.03472.
\bibitem{GGHV} J.~Guccione, J.~Guccione, C.~Horruitiner, J.~Valqui, \emph{Increasing
the degree bound}, arXiv:2204.14178.
\bibitem{LeeLi} K.~Lee, L.~Li, \emph{Magnus' formula revisited, IV},
arXiv:2408.01279.
\bibitem{Furter} J.-P.~Furter, \emph{On the variety of automorphisms of the affine
plane}, J.\ Algebra \textbf{195} (1997), 604--623.
\bibitem{Pinchuk} S.~Pinchuk, \emph{A counterexample to the strong real Jacobian
conjecture}, Math.\ Z.\ \textbf{217} (1994), 1--4.
\end{thebibliography}

\end{document}
